
\documentclass{article}
\usepackage{colortbl}
\usepackage{makecell}
\usepackage{multirow}
\usepackage{supertabular}

\begin{document}

\newcounter{utterance}

\centering \large Interaction Transcript for game `imagegame', experiment `random\_grids', episode 6 with Qwen3{-}30B{-}A3B{-}t0.0.
\vspace{24pt}

{ \footnotesize  \setcounter{utterance}{1}
\setlength{\tabcolsep}{0pt}
\begin{supertabular}{c@{$\;$}|p{.15\linewidth}@{}p{.15\linewidth}p{.15\linewidth}p{.15\linewidth}p{.15\linewidth}p{.15\linewidth}}
    \# & $\;$A & \multicolumn{4}{c}{Game Master} & $\;\:$B\\
    \hline

    \theutterance \stepcounter{utterance}  
    & & \multicolumn{4}{p{0.6\linewidth}}{
        \cellcolor[rgb]{0.9,0.9,0.9}{
            \makecell[{{p{\linewidth}}}]{
                \texttt{\tiny{[P1$\langle$GM]}}
                \texttt{Let us do the following task. The goal is to fill an empty grid that looks like this:} \\
\\ 
\texttt{▢ ▢ ▢ ▢ ▢} \\
\texttt{▢ ▢ ▢ ▢ ▢} \\
\texttt{▢ ▢ ▢ ▢ ▢} \\
\texttt{▢ ▢ ▢ ▢ ▢} \\
\texttt{▢ ▢ ▢ ▢ ▢} \\
\\ 
\texttt{A filled grid below is 5 by 5 and can look like this:} \\
\\ 
\texttt{▢ ▢ ▢ ▢ ▢} \\
\texttt{▢ ▢ E ▢ ▢} \\
\texttt{▢ ▢ ▢ ▢ ▢} \\
\texttt{▢ ▢ ▢ ▢ ▢} \\
\texttt{X X X X X} \\
\\ 
\texttt{I want you to describe this grid to me, step by step. You don't need to describe the empty squares, which are denoted with "▢". Only describe the location of letters in the grid. Then you wait for me to say "What is your next command?", and then you continue with the next step. Take the size of the grid into consideration while giving commands. You can give only one command at each turn but it can refer to single or multiple cells, rows, columns or even the whole grid at once. When you have described everything, you say "Command: DONE". Only give the command and omit everything else.} \\
\\ 
\texttt{For the filled grid above, here are the example steps.} \\
\\ 
\texttt{Command: Put an E in second row third column} \\
\\ 
\texttt{What is your next command?} \\
\texttt{Command: Fill the last row with X} \\
\\ 
\texttt{What is your next command?} \\
\texttt{Command: DONE} \\
\\ 
\texttt{Ok. Please do this for the following example, which is a 5 by 5 grid.} \\
\texttt{M ▢ ▢ ▢ ▢} \\
\texttt{M ▢ ▢ M ▢} \\
\texttt{▢ ▢ ▢ ▢ ▢} \\
\texttt{▢ ▢ M ▢ ▢} \\
\texttt{▢ M ▢ ▢ ▢} \\
\texttt{What is your next command?} \\
            }
        }
    }
    & \\ \\

    \theutterance \stepcounter{utterance}  
    & \multicolumn{4}{p{0.6\linewidth}}{
        \cellcolor[rgb]{0.9,0.9,0.9}{
            \makecell[{{p{\linewidth}}}]{
                \texttt{\tiny{[P1$\rangle$GM]}}
                \texttt{<think>} \\
\texttt{Okay, let's see. The user wants me to describe how to fill a 5x5 grid step by step, but only mention the cells that have letters. The example they gave was placing an E and then filling the last row with Xs. So I need to do the same for this new grid.} \\
\\ 
\texttt{First, let me look at the given grid:} \\
\\ 
\texttt{Row 1: M ▢ ▢ ▢ ▢} \\
\texttt{Row 2: M ▢ ▢ M ▢} \\
\texttt{Row 3: ▢ ▢ ▢ ▢ ▢} \\
\texttt{Row 4: ▢ ▢ M ▢ ▢} \\
\texttt{Row 5: ▢ M ▢ ▢ ▢} \\
\\ 
\texttt{So the letters are in specific positions. Let me note their coordinates. Remembering that rows are 1{-}5 and columns 1{-}5.} \\
\\ 
\texttt{Row 1, Column 1: M} \\
\texttt{Row 2, Column 1: M} \\
\texttt{Row 2, Column 4: M} \\
\texttt{Row 4, Column 3: M} \\
\texttt{Row 5, Column 2: M} \\
\\ 
\texttt{So there are five Ms. The task is to describe each of these positions step by step. But the user said to only describe the location of letters, not the empty squares. So each M needs to be placed.} \\
\\ 
\texttt{But how to do it step by step. The example had two steps: placing E and then filling a row. So maybe here, each M is} \\
            }
        }
    }
    & & \\ \\

    \theutterance \stepcounter{utterance}  
    & & & \multicolumn{2}{p{0.3\linewidth}}{
        \cellcolor[rgb]{0.9,0.9,0.9}{
            \makecell[{{p{\linewidth}}}]{
                \texttt{\tiny{[GM$|$GM]}}
                \texttt{Invalid instruction format} \\
            }
        }
    }
    & & \\ \\

\end{supertabular}
}

\end{document}
